%%
%% This is file `sample-sigconf.tex',
%% generated with the docstrip utility.
%%
%% The original source files were:
%%
%% samples.dtx  (with options: `all,proceedings,bibtex,sigconf')
%% 
%% IMPORTANT NOTICE:
%% 
%% For the copyright see the source file.
%% 
%% Any modified versions of this file must be renamed
%% with new filenames distinct from sample-sigconf.tex.
%% 
%% For distribution of the original source see the terms
%% for copying and modification in the file samples.dtx.
%% 
%% This generated file may be distributed as long as the
%% original source files, as listed above, are part of the
%% same distribution. (The sources need not necessarily be
%% in the same archive or directory.)
%%
%%
%% Commands for TeXCount
%TC:macro \cite [option:text,text]
%TC:macro \citep [option:text,text]
%TC:macro \citet [option:text,text]
%TC:envir table 0 1
%TC:envir table* 0 1
%TC:envir tabular [ignore] word
%TC:envir displaymath 0 word
%TC:envir math 0 word
%TC:envir comment 0 0
%%
%% The first command in your LaTeX source must be the \documentclass
%% command.
%%
%% For submission and review of your manuscript please change the
%% command to \documentclass[manuscript, screen, review]{acmart}.
%%
%% When submitting camera ready or to TAPS, please change the command
%% to \documentclass[sigconf]{acmart} or whichever template is required
%% for your publication.
%%
%%
\documentclass[sigconf]{acmart}
%%
%% \BibTeX command to typeset BibTeX logo in the docs
\AtBeginDocument{%
  \providecommand\BibTeX{{%
    Bib\TeX}}}

% remove "template" option to hide blue-font notes
\usepackage[template]{mom-template} 

\usepackage{multirow}
\usepackage{here}
\usepackage{xfp}
\makeatletter

\usepackage{xcolor}
\iftemplate
  \newcommand{\note}[1]{\textcolor{blue}{\textit{#1}}}
\else
  \newcommand{\note}[1]{}
\fi

% ============================================================
% CONFIGURATION
% ============================================================

% Total number of MTs
\def\MoM@mtcount{5}

% Number of questions per MT
\expandafter\def\csname MT@numq@1\endcsname{7}
\expandafter\def\csname MT@numq@2\endcsname{7}
\expandafter\def\csname MT@numq@3\endcsname{6}
\expandafter\def\csname MT@numq@4\endcsname{4}
\expandafter\def\csname MT@numq@5\endcsname{5}

% ============================================================
% NC / PC / FC  mappings
% ============================================================

% Numeric values (used for the radar average)
\expandafter\def\csname MT@val@NC\endcsname{0}
\expandafter\def\csname MT@val@PC\endcsname{0.5}
\expandafter\def\csname MT@val@FC\endcsname{1}

% Full text labels (used in the per-MT score table)
\expandafter\def\csname MT@text@NC\endcsname{Not covered}
\expandafter\def\csname MT@text@PC\endcsname{Partially covered}
\expandafter\def\csname MT@text@FC\endcsname{Fully covered}

% \MT@MapVal{NC|PC|FC} -> 0 / 0.5 / 1
\newcommand{\MT@MapVal}[1]{%
  \ifcsname MT@val@#1\endcsname
    \csname MT@val@#1\endcsname
  \else
    0%
  \fi
}

% \MT@MapText{NC|PC|FC} -> "Not covered" / "Partially covered" / "Fully covered"
\newcommand{\MT@MapText}[1]{%
  \ifcsname MT@text@#1\endcsname
    \csname MT@text@#1\endcsname
  \else
    \textit{(unset)}%
  \fi
}

% ============================================================
% INTERNALS 
% ============================================================

% \MTNumQ{n} -- expands to the number of questions for MT n
\newcommand{\MTNumQ}[1]{\csname MT@numq@#1\endcsname}

% \MT{<mt>}{<q>}{<coverage>}
% - if <coverage> is non-empty: store the symbolic value (NC/PC/FC)
% - if <coverage> is empty: retrieve it (undefined entries read as NC)
\newcommand{\MT}[3]{%
  \if\relax\detokenize{#3}\relax
    \ifcsname MT@#1@#2\endcsname
      \csname MT@#1@#2\endcsname
    \else
      NC%
    \fi
  \else
    \expandafter\gdef\csname MT@#1@#2\endcsname{#3}%
  \fi
}

% --- Loop counters (two to avoid conflicts in nested loops) ---
\newcount\MT@ctrA
\newcount\MT@ctrB

% ============================================================
% PER-MT AVERAGE (for radar)
% ============================================================

% \MT@BuildAvgExpr{mt} -- builds \MT@avgexpr as an fpeval-ready expression
%   that evaluates to the average mapped value for MT <mt>.
\newcommand{\MT@BuildAvgExpr}[1]{%
  \global\MT@ctrB=0\relax
  \gdef\MT@avgexpr{0}%
  \loop
    \global\advance\MT@ctrB by 1
    \ifnum\MT@ctrB<\numexpr\MTNumQ{#1}+1\relax
      \xdef\MT@avgexpr{%
        \unexpanded\expandafter{\MT@avgexpr}%
        +\noexpand\MT@MapVal{\noexpand\MT{#1}{\the\MT@ctrB}{}}}%
  \repeat
  % divide by number of questions
  \xdef\MT@avgexpr{(\unexpanded\expandafter{\MT@avgexpr})/\MTNumQ{#1}}%
}

% ============================================================
% TABLE ROW BUILDER (3 columns: MT | Question | Coverage)
% ============================================================

\newcommand{\MT@Row}[2]{%
  & \textbf{#2} & \textsf{\MT@MapText{\MT{#1}{#2}{}}} \\
}

\newcommand{\MT@BuildRows}[1]{%
  \global\MT@ctrB=1\relax
  \gdef\MT@builtrows{}%
  \loop
    \global\advance\MT@ctrB by 1
    \ifnum\MT@ctrB<\numexpr\MTNumQ{#1}+1\relax
      \xdef\MT@builtrows{%
        \unexpanded\expandafter{\MT@builtrows}%
        \noexpand\MT@Row{#1}{\the\MT@ctrB}}%
  \repeat
}

% Full multirow block for MT #1
\newcommand{\MT@Block}[1]{%
  \MT@BuildRows{#1}%
  \multirow{\MTNumQ{#1}}{4em}{\textsf{\textbf{MT.#1}}}%
  & \textbf{1} & \textsf{\MT@MapText{\MT{#1}{1}{}}} \\
  \MT@builtrows
  \hline
}

% ============================================================
% RADAR BUILDER
% ============================================================

% Build radar axes + coverage path covering ALL configured MTs
% (1..\MoM@mtcount), regardless of whether \MTScoreTable was called
% for them.  Omitted MTs simply contribute an average of 0 (NC).
%
% \MT@radaraxes -- TikZ \draw / \node commands for the spokes + labels
% \MT@radarpath -- sequence of  (<angle>:<r>) --  coordinates

% Per-MT step.  Lives outside the outer \loop so it can safely use its
% own \loop (via \MT@BuildAvgExpr) -- plain TeX's \loop\repeat is not
% nestable, so we collect a list of \MT@RadarStep calls during the
% outer loop and execute them afterwards.
\newcommand{\MT@RadarStep}[1]{%
  % angle = 90 - 360*(i-1)/N  (clockwise from top, first axis on top)
  \edef\MT@angle{\fpeval{90 - 360*(#1-1)/\MoM@mtcount}}%
  \MT@BuildAvgExpr{#1}%
  \edef\MT@avg{\fpeval{\MT@avgexpr}}%
  \xdef\MT@radaraxes{%
    \unexpanded\expandafter{\MT@radaraxes}%
    \noexpand\draw[gray!60] (0,0) -- (\MT@angle:1);
    \noexpand\node[font=\noexpand\small\noexpand\sffamily]
      at (\MT@angle:1.18) {MT.#1};
  }%
  \xdef\MT@radarpath{%
    \unexpanded\expandafter{\MT@radarpath}%
    (\MT@angle:\MT@avg) --
  }%
  \xdef\MT@radarframe{%
    \unexpanded\expandafter{\MT@radarframe}%
    (\MT@angle:1) --
  }%
}

\newcommand{\MT@BuildRadarPath}{%
  % Phase 1: queue one \MT@RadarStep call per MT (no nested \loop yet).
  \gdef\MT@radarcalls{}%
  \global\MT@ctrA=0\relax
  \loop
    \global\advance\MT@ctrA by 1
    \ifnum\MT@ctrA<\numexpr\MoM@mtcount+1\relax
      \xdef\MT@radarcalls{%
        \unexpanded\expandafter{\MT@radarcalls}%
        \noexpand\MT@RadarStep{\the\MT@ctrA}}%
  \repeat
  % Phase 2: execute the queued calls (each may safely use \loop).
  \gdef\MT@radaraxes{}%
  \gdef\MT@radarpath{}%
  \gdef\MT@radarframe{}%
  \MT@radarcalls
}

% ============================================================
% PUBLIC INTERFACE
% ============================================================

% \MTScoreTable{n} -- insert a coverage table for MT n
\newcommand{\MTScoreTable}[1]{%
  \medskip
  \begin{center}
  \begin{tabular}{ |c|c|c| }
  \cline{2-3}
  \multicolumn{1}{c|}{} & \textbf{Question} & \textbf{Coverage} \\
  \hline
  \MT@Block{#1}%
  \end{tabular}
  \end{center}
}

% \MTRadar -- draw a radar / spider diagram of per-MT average coverage
\newcommand{\MTRadar}{%
  \MT@BuildRadarPath
  \begin{center}
  \begin{tikzpicture}[scale=3]
    % Concentric reference polygons (NC=0, PC=0.5, FC=1 levels)
    \foreach \r in {0.25,0.5,0.75,1.0}{%
      \begin{scope}[scale=\r]
        \draw[gray!25, thin] \MT@radarframe cycle;
      \end{scope}%
    }
    % Spokes and MT labels
    \MT@radaraxes
    % Coverage polygon
    \draw[gray!60!black, thick, fill=gray!40, fill opacity=0.55]
      \MT@radarpath cycle;
  \end{tikzpicture}
  \end{center}
}

\makeatother

% MT1 (7 questions)
\MT{1}{1}{FC} 
\MT{1}{2}{NC} 
\MT{1}{3}{NC} 
\MT{1}{4}{PC} 
\MT{1}{5}{NC} 
\MT{1}{6}{FC} 
\MT{1}{7}{NC}

% MT2 (7 questions)
\MT{2}{1}{FC} 
\MT{2}{2}{FC} 
\MT{2}{3}{FC} 
\MT{2}{4}{FC} 
\MT{2}{5}{PC} 
\MT{2}{6}{FC} 
\MT{2}{7}{FC}

% MT3 (6 questions)
\MT{3}{1}{NC} 
\MT{3}{2}{NC} 
\MT{3}{3}{NC} 
\MT{3}{4}{FC} 
\MT{3}{5}{FC} 
\MT{3}{6}{NC}

% MT4 (4 questions)
\MT{4}{1}{NC} 
\MT{4}{2}{FC} 
\MT{4}{3}{FC} 
\MT{4}{4}{PC}

% MT5 (5 questions)
\MT{5}{1}{FC} 
\MT{5}{2}{NC} 
\MT{5}{3}{NC} 
\MT{5}{4}{PC} 
\MT{5}{5}{FC}


\usepackage{tikz}
\usepackage{xcolor}

%% Rights management information.  This information is sent to you
%% when you complete the rights form.  These commands have SAMPLE
%% values in them; it is your responsibility as an author to replace
%% the commands and values with those provided to you when you
%% complete the rights form.
\setcopyright{acmlicensed}
\copyrightyear{2026}
\acmYear{2026}
\acmDOI{XXXXXXX.XXXXXXX}
%% These commands are for a PROCEEDINGS abstract or paper.
\acmConference[MODELS '26]{Make sure to enter the correct
  conference title from your rights confirmation email}{June 03--05,
  2018}{Woodstock, NY}
%%
%%  Uncomment \acmBooktitle if the title of the proceedings is different
%%  from ``Proceedings of ...''!
%%
%%\acmBooktitle{Woodstock '18: ACM Symposium on Neural Gaze Detection,
%%  June 03--05, 2018, Woodstock, NY}
\acmISBN{978-1-4503-XXXX-X/2018/06}


%%
%% Submission ID.
%% Use this when submitting an article to a sponsored event. You'll
%% receive a unique submission ID from the organizers
%% of the event, and this ID should be used as the parameter to this command.
%%\acmSubmissionID{123-A56-BU3}

%%
%% For managing citations, it is recommended to use bibliography
%% files in BibTeX format.
%%
%% You can then either use BibTeX with the ACM-Reference-Format style,
%% or BibLaTeX with the acmnumeric or acmauthoryear sytles, that include
%% support for advanced citation of software artefact from the
%% biblatex-software package, also separately available on CTAN.
%%
%% Look at the sample-*-biblatex.tex files for templates showcasing
%% the biblatex styles.
%%

%%
%% The majority of ACM publications use numbered citations and
%% references.  The command \citestyle{authoryear} switches to the
%% "author year" style.
%%
%% If you are preparing content for an event
%% sponsored by ACM SIGGRAPH, you must use the "author year" style of
%% citations and references.
%% Uncommenting
%% the next command will enable that style.
%%\citestyle{acmauthoryear}


%%
%% end of the preamble, start of the body of the document source.
\begin{document}

%%
%% The "title" command has an optional parameter,
%% allowing the author to define a "short title" to be used in page headers.
\title{Template for Submission of Solution to the\\2026 Model Management (MoM) Challenge}

%%
%% The "author" command and its associated commands are used to define
%% the authors and their affiliations.
%% Of note is the shared affiliation of the first two authors, and the
%% "authornote" and "authornotemark" commands
%% used to denote shared contribution to the research.
\author{Jane Smith}
\authornote{Both authors contributed equally to this research.}
\email{jane@corporation.com}
\orcid{1234-5678-9012}
\author{John Smith}
\authornotemark[1]
\email{john@corporation.com}
\affiliation{%
  \institution{MoM Corporation}
  \city{Dublin}
  \state{Ohio}
  \country{USA}
}

\author{Jane Doe}
\affiliation{%
  \institution{Model Managers}
  \city{Hekla}
  \country{Iceland}}
\email{larst@affiliation.org}

%%
%% By default, the full list of authors will be used in the page
%% headers. Often, this list is too long, and will overlap
%% other information printed in the page headers. This command allows
%% the author to define a more concise list
%% of authors' names for this purpose.
\renewcommand{\shortauthors}{Smith, Smith \textit{et al.}}

%%
%% The abstract is a short summary of the work to be presented in the
%% article.
\begin{abstract}
  We contribute a solution to the Model Management (MoM) Challenge 2026 
(available at doi: 10.5281/zenodo.19608752) with the XYZ tool / framework. 

\note{Describe how your contribution addressed the challenge 
requirements on a high abstraction level.}

\note{Before submitting your manuscript remove the ``template'' option to hide blue-font notes from the preamble in \\\textbackslash{}usepackage[template]\{mom-template\}}

\note{\textbf{Please make sure to submit} your solution to the "Challenge" track on the MoM submission EasyChair site\\\url{https://easychair.org/conferences/?conf=mom2026}}
\end{abstract}

%%
%% The code below is generated by the tool at http://dl.acm.org/ccs.cfm.
%% Please copy and paste the code instead of the example below.
%%
\begin{CCSXML}
<ccs2012>
 <concept>
  <concept_id>00000000.0000000.0000000</concept_id>
  <concept_desc>Do Not Use This Code, Generate the Correct Terms for Your Paper</concept_desc>
  <concept_significance>500</concept_significance>
 </concept>
 <concept>
  <concept_id>00000000.00000000.00000000</concept_id>
  <concept_desc>Do Not Use This Code, Generate the Correct Terms for Your Paper</concept_desc>
  <concept_significance>100</concept_significance>
 </concept>
</ccs2012>
\end{CCSXML}

\ccsdesc[500]{Do Not Use This Code~Generate the Correct Terms for Your Paper}
\ccsdesc[300]{Do Not Use This Code~Generate the Correct Terms for Your Paper}

%%
%% Keywords. The author(s) should pick words that accurately describe
%% the work being presented. Separate the keywords with commas.
\keywords{Do, Not, Use, This, Code, Put, the, Correct, Terms, for,
  Your, Paper}

%\received{20 February 2007}
%\received[revised]{12 March 2009}
%\received[accepted]{5 June 2009}

%%
%% This command processes the author and affiliation and title
%% information and builds the first part of the formatted document.
\maketitle

\section{Introduction}
\label{sec:Introduction}

\note{
General introduction of the solution:
\begin{itemize}
    \item Why did you chose to take on the Challenge?
    \item What are the general capabilities of your tool?
    \item Since when is the tool being developed?
    \item Is it open-source? 
    \item Is it dependent on another platform (Eclipse, etc.)
    \item Is it a commercial tool? Who are the customers?
\end{itemize}
}

\section{High-Level Dimensions and Functions of MoM}
\label{sec:HDF}

    \note{Position your tool in the context of high-level dimensions 
    and functions of MoM as described in the Challenge.}
    
\subsection{Approach [HDF.1]}
\label{subsec:HDF-1}

\subsection{Technological Space [HDF.2]}
\label{subsec:HDF-2}

\subsection{Relationships [HDF.3]}
\label{subsec:HDF-3}

\subsection{Views [HDF.4]}
\label{subsec:HDF-4}

\subsection{Collaboration [HDF.5]}
\label{subsec:HDF-5}

\section{Model Management Tasks}
\label{sec:MT}

    \note{It is perfectly acceptable if your solution only covers some and not all of the MoM Tasks.}
    
    \note{To record coverage for each MT question, open \texttt{score-tables.tex}
    and set the third argument of each \texttt{\textbackslash{}MT} macro to one of:
    \texttt{NC} (Not covered), \texttt{PC} (Partially covered), or \texttt{FC} (Fully covered).
    For example, marking \textsf{MT.2.3} as fully covered translates to
    \texttt{\textbackslash{}MT\{2\}\{3\}\{FC\}}.}
    
    \note{If an MT is not addressed by your solution, leave its questions as \texttt{NC} in \texttt{score-tables.tex} (the default). It will still appear as an axis on the coverage radar diagram, just at a value of zero. Make sure to mention something like ``This task is not handled by our tool.'' in the relevant section.}
    

%%% --------------------------------------------------------
%%% MT.1  Change Impact & Reconciliation
%%% --------------------------------------------------------
\subsection{Change Impact \& Reconciliation [MT.1]}
\label{sec:MT1}

\subsubsection*{MT.1.1 --- How is the change in the Parts Catalogue detected?}
\label{sec:MT1-1}

    \note{FC if your tool can detect and report changes; NC otherwise.}


\subsubsection*{MT.1.2 --- Which mechanisms or processes ensure propagation of this change to the Bill of Materials?}
\label{sec:MT1-2}

    \note{FC if a propagation mechanism exists; NC otherwise.}


\subsubsection*{MT.1.3 --- Is the propagation automatic, semi-automatic, or manual?}
\label{sec:MT1-3}

    \note{FC if propagation is at least semi-automatic; PC if manual but documented; NC if unsupported.}


\subsubsection*{MT.1.4 --- Can the propagation strategy be parameterised (to happen semi-/automatically)?}
\label{sec:MT1-4}

    \note{FC if the strategy is configurable; NC otherwise.}


\subsubsection*{MT.1.5 --- At which granularity does propagation operate (e.g., element, feature, model)?}
\label{sec:MT1-5}

    \note{FC if granularity can be described; NC if propagation is unsupported.}


\subsubsection*{MT.1.6 --- Can an explicit structure (e.g., graph, tree) be computed to represent impacted elements?}
\label{sec:MT1-6}

    \note{FC if an impact structure can be produced; NC otherwise.}


\subsubsection*{MT.1.7 --- Can such structures be composed across multiple changes, and how deeply can they be inspected?}
\label{sec:MT1-7}

    \note{FC if composition and inspection are supported; NC otherwise.}


\subsubsection*{Alternatives \& Extensions}
\label{sec:MT1-Alternatives-Extensions}

    \note{Discuss alternative mechanisms or extensions beyond the 
    scope of the questions above, if any.}

\subsubsection*{Score}
\label{sec:MT1-Score}

\MTScoreTable{1}


%%% --------------------------------------------------------
%%% MT.2  Views Management
%%% --------------------------------------------------------
\subsection{Views Management [MT.2]}
\label{sec:MT2}

\subsubsection*{MT.2.1 --- Can model elements or types be selectively filtered on demand? On which criteria?}
\label{sec:MT2-1}

\note{FC if selective filtering is supported; NC otherwise.}


\subsubsection*{MT.2.2 --- Can filtering rules be specified declaratively (e.g., via a query language)?}
\label{sec:MT2-2}

    \note{FC if declarative specification is supported; NC otherwise.}


\subsubsection*{MT.2.3 --- Can views be systematically derived from conformance models? How?}
\label{sec:MT2-3}

    \note{FC if systematic derivation is supported; NC otherwise.}


\subsubsection*{MT.2.4 --- Are views themselves typed (e.g., via a metamodel)?}
\label{sec:MT2-4}

\note{FC if views are typed; NC otherwise.}


\subsubsection*{MT.2.5 --- How are changes in views reconciled with the source model?}
\label{sec:MT2-5}

\note{FC if view-to-source reconciliation is supported; NC otherwise.}


\subsubsection*{MT.2.6 --- How are views notified of changes in the source model?}
\label{sec:MT2-6}

\note{FC if a notification or synchronisation mechanism exists; NC otherwise.}


\subsubsection*{MT.2.7 --- Do source model changes automatically propagate to views? How?}
\label{sec:MT2-7}

\note{FC if automatic propagation is supported; NC otherwise.}


\subsubsection*{Alternatives \& Extensions}
\label{sec:MT2-Alternatives-Extensions}

\note{Discuss alternative mechanisms or extensions beyond the 
scope of the questions above, if any.}

\subsubsection*{Score}
\label{sec:MT2-Score}

\MTScoreTable{2}


%%% --------------------------------------------------------
%%% MT.3  Querying & Validation
%%% --------------------------------------------------------
\subsection{Querying \& Validation [MT.3]}
\label{sec:MT3}

\subsubsection*{MT.3.1 --- Is \textsf{calculated\_total\_mass\_kg} automatically updated after the change?}
\label{sec:MT3-1}

\note{FC if automatic update is supported; NC otherwise.}


\subsubsection*{MT.3.2 --- Which mechanisms or processes support this update?}
\label{sec:MT3-2}

\note{FC if a concrete mechanism is described; NC otherwise.}


\subsubsection*{MT.3.3 --- Does it require manual intervention?}
\label{sec:MT3-3}

\note{FC if no manual intervention is needed; PC if partially manual; NC if fully manual or unsupported.}


\subsubsection*{MT.3.4 --- Can the process be fully automated or parameterised?}
\label{sec:MT3-4}

\note{FC if automation is configurable; NC otherwise.}


\subsubsection*{MT.3.5 --- Is the updated status reflected in the Report?}
\label{sec:MT3-5}

\note{FC if the Report is updated accordingly; NC otherwise.}


\subsubsection*{MT.3.6 --- Can impact analysis be performed manually by querying all relevant models?}
\label{sec:MT3-6}

\note{FC if cross-model querying is supported; NC otherwise.}


\subsubsection*{Alternatives \& Extensions}
\label{sec:MT3-Alternatives-Extensions}

\note{Discuss alternative mechanisms or extensions beyond the 
scope of the questions above, if any.}

\subsubsection*{Score}
\label{sec:MT3-Score}

\MTScoreTable{3}


%%% --------------------------------------------------------
%%% MT.4  Concurrent Modifications
%%% --------------------------------------------------------
\subsection{Concurrent Modifications [MT.4]}
\label{sec:MT4}

\subsubsection*{MT.4.1 --- How are conflicts resolved: manual vs.\ (semi-)automatic, online vs.\ offline, configurable vs.\ fixed strategies?}
\label{sec:MT4-1}

\note{FC if a conflict resolution strategy is supported; NC otherwise.}


\subsubsection*{MT.4.2 --- Are changes propagated in real time or through explicit synchronization?}
\label{sec:MT4-2}

\note{FC if propagation mode is described; NC if concurrent editing is unsupported.}


\subsubsection*{MT.4.3 --- If concurrent conflicting modifications are prevented, which mechanisms or policies detect and avoid such conflicts?}
\label{sec:MT4-3}

\note{FC if prevention mechanisms are described; NC if neither prevention nor resolution is supported.}


\subsubsection*{MT.4.4 --- Does the tool support live collaborative modelling? If so, how are conflicts visualized, prevented, or resolved?}
\label{sec:MT4-4}

\note{FC if live collaboration is supported; NC otherwise.}


\subsubsection*{Alternatives \& Extensions}
\label{sec:MT4-Alternatives-Extensions}

\note{Discuss alternative mechanisms or extensions beyond the 
scope of the questions above, if any.}

\subsubsection*{Score}
\label{sec:MT4-Score}

\MTScoreTable{4}


%%% --------------------------------------------------------
%%% MT.5  Version Management
%%% --------------------------------------------------------
\subsection{Version Management [MT.5]}
\label{sec:MT5}

\subsubsection*{MT.5.1 --- What is the definition of a ``version'' (single artefact, dependency closure, or system-wide state)?}
\label{sec:MT5-1}

\note{FC if the notion of version is clearly defined; NC otherwise.}


\subsubsection*{MT.5.2 --- Are versions created automatically or explicitly (e.g., commit-based)?}
\label{sec:MT5-2}

\note{FC if version creation is supported; NC otherwise.}


\subsubsection*{MT.5.3 --- How are artefacts linked across versions (e.g., traceability links between instances)?}
\label{sec:MT5-3}

\note{FC if cross-version linking is supported; NC otherwise.}


\subsubsection*{MT.5.4 --- Can versions be enriched with metadata?}
\label{sec:MT5-4}

\note{FC if metadata attachment is supported; NC otherwise.}


\subsubsection*{MT.5.5 --- Are such metadata queryable and exploitable?}
\label{sec:MT5-5}

\note{FC if metadata can be queried; NC otherwise.}


\subsubsection*{MT.5.6 --- Can the rationale for version creation be captured and maintained?}
\label{sec:MT5-6}

\note{FC if rationale capture is supported; NC otherwise.}


\subsubsection*{Alternatives \& Extensions}
\label{sec:MT5-Alternatives-Extensions}

\note{Discuss alternative mechanisms or extensions beyond the 
scope of the questions above, if any.}

\subsubsection*{Score}
\label{sec:MT5-Score}

\MTScoreTable{5}

\section{Replication Package (Optional)}
\label{sec:RP}

\note{We recommend authors provide a replication package to accompany their solution. Simple impermanent links to the source code are not acceptable. Consider using technologies like containerization, notebooks, etc. in order to enhance the replicability of your solution, and deposit them on platforms with some kind of guarantee of permanence like Zenodo. Ensure you provide complete instructions for not just initializing your tool, but also to reproduce the Engineering Scenarios (ES) described in the Challenge.}

\section{Conclusion}
\label{sec:Conclusion}

\note{
\begin{itemize}
   \item Summarize the \textbf{main strengths and limitations} of the approach;
   \item Describe any \textbf{adaptations or assumptions} made to fit the Case Study;
   \item Identify \textbf{unaddressed requirements or limitations} of the challenge;
   \item Provide \textbf{lessons learnt} and implications for future work,
   both for the authors and the broader MoM community.
\end{itemize}
}

%%
%% The acknowledgments section is defined using the "acks" environment
%% (and NOT an unnumbered section). This ensures the proper
%% identification of the section in the article metadata, and the
%% consistent spelling of the heading.
\begin{acks}
To XX for YY.
\end{acks}

%%
%% The next two lines define the bibliography style to be used, and
%% the bibliography file.
\bibliographystyle{ACM-Reference-Format}
\bibliography{sample-base}


%%
%% If your work has an appendix, this is the place to put it.
\appendix

\section{Coverage Summary}
\label{app:Scores}

\MTRadar

\end{document}
\endinput
%%
%% End of file `sample-sigconf.tex'.